\documentclass[11pt]{article}

\usepackage[margin=1in]{geometry}
\usepackage{amsmath, amssymb, amsthm}
\usepackage{mathtools}
\usepackage{parskip}
\usepackage{hyperref}

\title{Toward the Cs\'asz\'ar Polyhedron via Cell-Complex Enumeration\\[2pt]
\large Step 1: Notation and the First Conditions}
\author{}
\date{}

\begin{document}
\maketitle

\section*{Notation}

Let $P_i = (x_i, y_i, z_i) \in \mathbb{R}^3$ be the coordinate vector of the symbolic
point $i$. Define the \emph{bracket}
\[
[abcd] \;:=\; \det\!\begin{pmatrix} 1 & x_a & y_a & z_a \\ 1 & x_b & y_b & z_b \\ 1 & x_c & y_c & z_c \\ 1 & x_d & y_d & z_d \end{pmatrix}
\;=\; \det\!\bigl(P_b-P_a,\; P_c-P_a,\; P_d-P_a\bigr),
\]
which is $6$ times the signed volume of the tetrahedron $abcd$ and is alternating in
$a,b,c,d$. With this,
\[
abc^{\pm} \;:=\; \bigl\{\,p\in\mathbb{R}^3 \;:\; \pm[a,b,c,p] > 0\,\bigr\}.
\]

\section*{Condition 1: \texorpdfstring{$3 \in 012^+$}{3 in 012+}}

This is just $[0,1,2,3] > 0$. Three equivalent forms:

\paragraph{(a) Determinant.}
\[
\det\!\begin{pmatrix} 1 & x_0 & y_0 & z_0 \\ 1 & x_1 & y_1 & z_1 \\ 1 & x_2 & y_2 & z_2 \\ 1 & x_3 & y_3 & z_3 \end{pmatrix} > 0
\quad\Longleftrightarrow\quad
\det\!\begin{pmatrix} x_1-x_0 & y_1-y_0 & z_1-z_0 \\ x_2-x_0 & y_2-y_0 & z_2-z_0 \\ x_3-x_0 & y_3-y_0 & z_3-z_0 \end{pmatrix} > 0
\]

\paragraph{(b) Scalar triple product (cross-product antisymmetry).}
\[
(P_1 - P_0)\,\cdot\,\bigl((P_2 - P_0)\times(P_3 - P_0)\bigr) \;>\; 0.
\]
This form privileges $P_0$ as a basepoint, hiding the $S_4$-antisymmetry on the
four points. \emph{Fully unparenthesized.} Distribute the differences (using
$a\!\times\!a=0$ and the cyclic identity $a\!\cdot\!(b\!\times\!c)=c\!\cdot\!(a\!\times\!b)$);
the bracket reorganizes into the alternating sum over $3$-element subsets of
$\{0,1,2,3\}$:
\[
[0,1,2,3] \;=\; P_1\!\cdot\!(P_2\!\times\!P_3) \;-\; P_0\!\cdot\!(P_2\!\times\!P_3) \;+\; P_0\!\cdot\!(P_1\!\times\!P_3) \;-\; P_0\!\cdot\!(P_1\!\times\!P_2).
\]
Equivalently, with $T(i,j,k) := P_i\cdot(P_j\!\times\!P_k)$ (cyclically symmetric in
$i,j,k$) and $\widehat{d}$ denoting the elements of $\{0,1,2,3\}\setminus\{d\}$ in
increasing order,
\[
[0,1,2,3] \;=\; \sum_{d=0}^{3} (-1)^d\, T\!\bigl(\widehat{d}\bigr).
\]
The $S_4$-antisymmetry is now manifest --- any swap of two points flips the
bracket's sign.

\paragraph{(c) Polynomial in coordinates.}
With $\Delta_i := P_i - P_0$ for $i=1,2,3$:
\[
\Delta_{1,x}(\Delta_{2,y}\Delta_{3,z} - \Delta_{2,z}\Delta_{3,y})
+ \Delta_{1,y}(\Delta_{2,z}\Delta_{3,x} - \Delta_{2,x}\Delta_{3,z})
+ \Delta_{1,z}(\Delta_{2,x}\Delta_{3,y} - \Delta_{2,y}\Delta_{3,x}) \;>\; 0.
\]
Expanding the $\Delta$'s as differences of original coordinates:
\[
\begin{aligned}
&(x_1{-}x_0)\bigl[(y_2{-}y_0)(z_3{-}z_0) - (z_2{-}z_0)(y_3{-}y_0)\bigr] \\
+\;&(y_1{-}y_0)\bigl[(z_2{-}z_0)(x_3{-}x_0) - (x_2{-}x_0)(z_3{-}z_0)\bigr] \\
+\;&(z_1{-}z_0)\bigl[(x_2{-}x_0)(y_3{-}y_0) - (y_2{-}y_0)(x_3{-}x_0)\bigr] \;>\; 0.
\end{aligned}
\]
Again the $S_4$-antisymmetry is hidden inside the differences.
\emph{Fully unparenthesized.} Distributing all subtractions yields the $24$-term
Levi-Civita expansion:
\[
[0,1,2,3] \;=\; \sum_{(d,a,b,c)\in S_4} \mathrm{sgn}(d,a,b,c)\, x_a\, y_b\, z_c,
\]
where $(d,a,b,c)$ ranges over permutations of $(0,1,2,3)$ and $d$ is the row
contributing the constant $1$ from the bordered determinant (so $d$ enters only
through the sign). Grouped by missing index $d$:
\[
\begin{aligned}
d{=}0:\quad &+\,x_1 y_2 z_3 - x_1 y_3 z_2 - x_2 y_1 z_3 + x_2 y_3 z_1 + x_3 y_1 z_2 - x_3 y_2 z_1 \\
d{=}1:\quad &-\,x_0 y_2 z_3 + x_0 y_3 z_2 + x_2 y_0 z_3 - x_2 y_3 z_0 - x_3 y_0 z_2 + x_3 y_2 z_0 \\
d{=}2:\quad &+\,x_0 y_1 z_3 - x_0 y_3 z_1 - x_1 y_0 z_3 + x_1 y_3 z_0 + x_3 y_0 z_1 - x_3 y_1 z_0 \\
d{=}3:\quad &-\,x_0 y_1 z_2 + x_0 y_2 z_1 + x_1 y_0 z_2 - x_1 y_2 z_0 - x_2 y_0 z_1 + x_2 y_1 z_0
\end{aligned}
\]
Every monomial uses each index at most once, and any swap of two indices flips the
signs of all $24$ terms in lock-step --- the four-fold antisymmetry on $\{0,1,2,3\}$
is entirely visible.

\section*{Condition 2: \texorpdfstring{$012^+ \cap 013^+ \neq \emptyset$}{012+ cap 013+ nonempty}}

Existence of $p$ with $[0,1,2,p] > 0$ \emph{and} $[0,1,3,p] > 0$.

\paragraph{(a) Determinant.}
\[
\exists\, p\in\mathbb{R}^3:\quad
\det\!\begin{pmatrix} 1 & x_0 & y_0 & z_0 \\ 1 & x_1 & y_1 & z_1 \\ 1 & x_2 & y_2 & z_2 \\ 1 & p_x & p_y & p_z \end{pmatrix} > 0
\;\wedge\;
\det\!\begin{pmatrix} 1 & x_0 & y_0 & z_0 \\ 1 & x_1 & y_1 & z_1 \\ 1 & x_3 & y_3 & z_3 \\ 1 & p_x & p_y & p_z \end{pmatrix} > 0
\]

\paragraph{(b) Scalar triple product.}
\[
\exists\, p\in\mathbb{R}^3:\quad
(P_1{-}P_0)\cdot\bigl((P_2{-}P_0)\times(p{-}P_0)\bigr) > 0
\;\wedge\;
(P_1{-}P_0)\cdot\bigl((P_3{-}P_0)\times(p{-}P_0)\bigr) > 0.
\]
\emph{Fully unparenthesized.} Each bracket cofactor-expands as in
Condition~1(b), with $p$ playing the role of the fourth point:
\[
\begin{aligned}
{[0,1,2,p]} \;&=\; P_1\!\cdot\!(P_2\!\times\!p) \;-\; P_0\!\cdot\!(P_2\!\times\!p) \;+\; P_0\!\cdot\!(P_1\!\times\!p) \;-\; P_0\!\cdot\!(P_1\!\times\!P_2),\\
{[0,1,3,p]} \;&=\; P_1\!\cdot\!(P_3\!\times\!p) \;-\; P_0\!\cdot\!(P_3\!\times\!p) \;+\; P_0\!\cdot\!(P_1\!\times\!p) \;-\; P_0\!\cdot\!(P_1\!\times\!P_3).
\end{aligned}
\]
The shared term $P_0\!\cdot\!(P_1\!\times\!p)$ reflects that both brackets share
the edge $\overline{P_0P_1}$. Each is $S_4$-antisymmetric in its own four points
($\{0,1,2,p\}$ resp.\ $\{0,1,3,p\}$) and \emph{affine} in $p$.

\paragraph{(c) Polynomial in coordinates.}
Substitute $p$ for $P_3$ in the first inequality, and $p$ for $P_2$ (keeping $P_3$
in the third slot) in the second. Writing $p=(x_p, y_p, z_p)$:
\[
\exists\,(x_p, y_p, z_p)\in\mathbb{R}^3:
\]
\[
\begin{aligned}
&(x_1{-}x_0)\bigl[(y_2{-}y_0)(z_p{-}z_0) - (z_2{-}z_0)(y_p{-}y_0)\bigr] \\
+\;&(y_1{-}y_0)\bigl[(z_2{-}z_0)(x_p{-}x_0) - (x_2{-}x_0)(z_p{-}z_0)\bigr] \\
+\;&(z_1{-}z_0)\bigl[(x_2{-}x_0)(y_p{-}y_0) - (y_2{-}y_0)(x_p{-}x_0)\bigr] \;>\; 0
\end{aligned}
\]
\[
\text{AND}
\]
\[
\begin{aligned}
&(x_1{-}x_0)\bigl[(y_3{-}y_0)(z_p{-}z_0) - (z_3{-}z_0)(y_p{-}y_0)\bigr] \\
+\;&(y_1{-}y_0)\bigl[(z_3{-}z_0)(x_p{-}x_0) - (x_3{-}x_0)(z_p{-}z_0)\bigr] \\
+\;&(z_1{-}z_0)\bigl[(x_3{-}x_0)(y_p{-}y_0) - (y_3{-}y_0)(x_p{-}x_0)\bigr] \;>\; 0.
\end{aligned}
\]
\emph{Fully unparenthesized.} Each bracket admits the $24$-term Levi-Civita
expansion of Condition~1(c). For $[0,1,2,p]$, treating $p$ as the fourth point
and grouping by missing index:
\[
\begin{aligned}
d{=}0:\quad &+\,x_1 y_2 z_p - x_1 y_p z_2 - x_2 y_1 z_p + x_2 y_p z_1 + x_p y_1 z_2 - x_p y_2 z_1 \\
d{=}1:\quad &-\,x_0 y_2 z_p + x_0 y_p z_2 + x_2 y_0 z_p - x_2 y_p z_0 - x_p y_0 z_2 + x_p y_2 z_0 \\
d{=}2:\quad &+\,x_0 y_1 z_p - x_0 y_p z_1 - x_1 y_0 z_p + x_1 y_p z_0 + x_p y_0 z_1 - x_p y_1 z_0 \\
d{=}p:\quad &-\,x_0 y_1 z_2 + x_0 y_2 z_1 + x_1 y_0 z_2 - x_1 y_2 z_0 - x_2 y_0 z_1 + x_2 y_1 z_0.
\end{aligned}
\]
$[0,1,3,p]$ is obtained by replacing every subscript $2$ with $3$ above. Two
structural observations are now visible:
\begin{itemize}
  \item Rows $d=0,1,2$ contain all the $p$-dependence (linear in $p$); row $d=p$
  is the constant $-\det(P_0,P_1,P_2)$ (resp.\ $-\det(P_0,P_1,P_3)$). So each
  bracket is affine in $p$, and the realizability question is whether two open
  halfspaces in $p$-space meet.
  \item The $S_4$-antisymmetry on the bracket's four points $\{0,1,2,p\}$ (or
  $\{0,1,3,p\}$) is manifest: every swap of two indices flips the signs of all
  $24$ monomials simultaneously.
\end{itemize}

\section*{Realizability: \texorpdfstring{$012^+\cap 013^+$}{012+ cap 013+} exists, \texorpdfstring{$012^-\cap 013^+\cap 123^+\cap 203^+$}{012- cap 013+ cap 123+ cap 203+} does not}

\paragraph{Notation conversion.} The bracket $[a,b,c,d]$ is alternating, so
$[2,0,3,p] = -[0,2,3,p]$, hence
\[
203^+ \;=\; \{p : [2,0,3,p] > 0\} \;=\; \{p : [0,2,3,p] < 0\} \;=\; 023^-.
\]
The cell in question is therefore $012^- \cap 013^+ \cap 023^- \cap 123^+$ in
canonical (sorted-index) notation --- the sign vector $(-,+,-,+)$ on the four
faces $(012, 013, 023, 123)$ of tetrahedron $0123$.

\paragraph{Key identity.} For any $p \in \mathbb{R}^3$,
\begin{equation}\label{eq:star}\tag{$\star$}
[0,1,2,p] \;-\; [0,1,3,p] \;+\; [0,2,3,p] \;-\; [1,2,3,p] \;=\; [0,1,2,3].
\end{equation}

\textit{Proof.} Since $P_0, P_1, P_2, P_3$ are affinely independent (because
$[0,1,2,3] \neq 0$), every $p \in \mathbb{R}^3$ has unique barycentric coordinates
\[
p \;=\; \alpha_0 P_0 + \alpha_1 P_1 + \alpha_2 P_2 + \alpha_3 P_3,
\qquad \alpha_0+\alpha_1+\alpha_2+\alpha_3 \;=\; 1.
\]
The lifted equality $(1,p) = \sum_i \alpha_i\,(1,P_i)$ holds in $\mathbb{R}^4$, so
multilinearity of the bordered $4\times 4$ determinant in row $p$ gives, for
distinct $a,b,c \in \{0,1,2,3\}$ with missing index $d$,
\[
[a,b,c,p] \;=\; \sum_{i=0}^{3} \alpha_i\, [a,b,c,i] \;=\; \alpha_d\, [a,b,c,d],
\]
since $[a,b,c,a] = [a,b,c,b] = [a,b,c,c] = 0$ (repeated row). Using
$[0,1,3,2] = -[0,1,2,3]$, $[0,2,3,1] = +[0,1,2,3]$, $[1,2,3,0] = -[0,1,2,3]$
(reorder the missing index back into canonical position), and dividing by
$[0,1,2,3]$:
\[
\alpha_3 \;=\; \frac{[0,1,2,p]}{[0,1,2,3]},\quad
\alpha_2 \;=\; \frac{-[0,1,3,p]}{[0,1,2,3]},\quad
\alpha_1 \;=\; \frac{[0,2,3,p]}{[0,1,2,3]},\quad
\alpha_0 \;=\; \frac{-[1,2,3,p]}{[0,1,2,3]}.
\]
Then $\alpha_0 + \alpha_1 + \alpha_2 + \alpha_3 = 1$, multiplied through by
$[0,1,2,3]$, is exactly~\eqref{eq:star}. \hfill$\square$

\paragraph{Why \texorpdfstring{$012^+ \cap 013^+$}{012+ cap 013+} is non-empty.}
Take the explicit witness
\[
p^\ast \;:=\; 2 P_3 - P_2 \qquad \text{(the reflection of } P_2 \text{ across } P_3\text{).}
\]
Its coefficients $(2, -1)$ sum to $1$, so any affine function of $p$ evaluated at
$p^\ast$ equals $2$ times its value at $P_3$ minus its value at $P_2$. With
$[0,1,2,P_2]=0$ and $[0,1,2,P_3]=[0,1,2,3]$:
\[
[0,1,2,\,p^\ast] \;=\; 2\,[0,1,2,3] \;-\; 0 \;=\; 2\,[0,1,2,3] \;>\; 0.
\]
With $[0,1,3,P_2]=[0,1,3,2]=-[0,1,2,3]$ and $[0,1,3,P_3]=0$:
\[
[0,1,3,\,p^\ast] \;=\; 2\cdot 0 \;-\; (-[0,1,2,3]) \;=\; [0,1,2,3] \;>\; 0.
\]
Both brackets are positive, so $p^\ast \in 012^+ \cap 013^+$. (Geometrically:
planes $012$ and $013$ are distinct and share the line $\overline{P_0 P_1}$, so
they cut $\mathbb{R}^3$ into four open wedges, all non-empty; $012^+\cap 013^+$
is one of them.)

\paragraph{Why \texorpdfstring{$012^- \cap 013^+ \cap 023^- \cap 123^+$}{012- cap 013+ cap 023- cap 123+} is empty.}
Suppose for contradiction such a $p$ exists. The four memberships translate to
\[
[0,1,2,p] < 0,\quad [0,1,3,p] > 0,\quad [0,2,3,p] < 0,\quad [1,2,3,p] > 0.
\]
Hence each summand on the left side of~\eqref{eq:star}, with its prescribed sign,
is \emph{strictly negative}:
\[
\underbrace{[0,1,2,p]}_{<\,0}
\;\underbrace{-\,[0,1,3,p]}_{<\,0}
\;\underbrace{+\,[0,2,3,p]}_{<\,0}
\;\underbrace{-\,[1,2,3,p]}_{<\,0}
\;<\; 0.
\]
But by~\eqref{eq:star} this sum equals $[0,1,2,3]$, which is $>0$ by Condition~1.
Contradiction. Hence $012^- \cap 013^+ \cap 023^- \cap 123^+ = \emptyset$.
\hfill$\square$

\paragraph{Geometric meaning.} In barycentric coordinates the four sign
conditions read
\[
\alpha_0 < 0,\quad \alpha_1 < 0,\quad \alpha_2 < 0,\quad \alpha_3 < 0
\qquad\text{(``$p$ on the wrong side of every face'')},
\]
which is impossible because $\alpha_0+\alpha_1+\alpha_2+\alpha_3 = 1 > 0$. The
four face-planes of $0123$ in general position cut $\mathbb{R}^3$ into $15$ open
cells (the count $\binom{4}{0}+\binom{4}{1}+\binom{4}{2}+\binom{4}{3}$), and the
missing $16^{\mathrm{th}}$ sign vector is precisely $(012^-, 013^+, 023^-, 123^+)$.
The asymmetry between the two cells of the question is now visible: $012^+\cap 013^+$
slices $\mathbb{R}^3$ along only \emph{two} planes (a wedge --- always non-empty),
while the missing $16^{\mathrm{th}}$ cell asks for a sign-coherent intersection
across \emph{all four} faces of a closed tetrahedron, which the affine relation
$\sum\alpha_i = 1$ rules out.

\section*{Observations}

\begin{itemize}
  \item \textbf{Only signs matter.} The cell complex depends only on the
  \emph{sign vector} in $\{+,-\}^{\binom{n}{3}}$. Coordinates are useful only as a
  realizability oracle: if we want a ``without coordinates'' enumeration, the
  coordinate forms are witnesses we consult only when the combinatorial sign rules
  cannot decide.

  \item \textbf{Condition 2 is automatic.} Planes $012$ and $013$ share the line
  $\overline{P_0P_1}$ and (since $P_3 \notin$ plane $012$ by Condition~1) are
  distinct. They partition $\mathbb{R}^3\setminus\overline{P_0P_1}$ into four open
  wedges, all non-empty. So $012^+\cap 013^+\neq\emptyset$ trivially. A concrete
  witness: any point $P_3 + \varepsilon\,\hat n_{013^+}$ with $\varepsilon$ small.

  \item \textbf{Where it gets interesting.} Once we have all four planes
  $\{012,\,013,\,023,\,123\}$ they cut $\mathbb{R}^3$ into
  $\binom{4}{0}+\binom{4}{1}+\binom{4}{2}+\binom{4}{3}=15$ regions. Of the $2^4=16$
  a-priori sign vectors, exactly one is empty --- the missing 16th cell. A clean
  way to spot it: the sign vector that says ``$p$ is on the \emph{opposite} side of
  every face from the opposite vertex of tetrahedron $0123$'' is not realizable
  (no point sees every face from the wrong side).
\end{itemize}

\section*{Algorithm: deciding if a new plane $ij(n{+}1)$ slices an existing cell}

We are now ready for the inductive step. Assume we have introduced $n+1$ named
points $\{P_0,\ldots,P_n\}$ and constructed the cell complex of the
$\binom{n+1}{3}$ planes through every triple. We pick a new point $P_{n+1}$
inside one specific cell $C^\ast$ of the current complex; this pins the signs of
all $4$-brackets involving $n{+}1$ (because membership of $P_{n+1}$ in $C^\ast$
is exactly $\sigma_{C^\ast}(\Pi) = \mathrm{sgn}([\Pi, n{+}1])$ for every existing
plane $\Pi$). Then we must slice every existing cell with the $\binom{n+1}{2}$
new planes $\bigl\{\,ij(n{+}1) : 0\le i<j\le n\,\bigr\}$. The basic question is:
given a cell $C$ and a new plane $\Pi = ij(n{+}1)$, does $\Pi$ slice $C$ (i.e.\
are both $C\cap\Pi^+$ and $C\cap\Pi^-$ non-empty)?

The algorithm below answers this question using the chirotope alone (no
coordinates), and is correct in the general case. The conceptual primitive is
the affine identity \eqref{eq:star} from earlier --- now applied to \emph{every}
$4$-tuple of named points, not just $\{0,1,2,3\}$.

\paragraph{State.} We maintain
\begin{itemize}
  \item the \textbf{chirotope}: $\mathrm{sgn}[a,b,c,d]$ for every $4$-subset
  $\{a,b,c,d\}\subseteq\{0,\ldots,n{+}1\}$;
  \item the \textbf{cell complex} of the previous step: each open cell $C$
  stored as a sign vector $\sigma_C : \mathcal{H}_{\mathrm{old}} \to \{+,-\}$
  over the old planes $\mathcal{H}_{\mathrm{old}} = \binom{\{0,\ldots,n\}}{3}$.
\end{itemize}
We do \emph{not} need anything else --- no coordinates, no magnitudes, no
``rays at infinity.''

\paragraph{The single sign-rule.} For every $4$-subset $\{a,b,c,d\}$ of named
indices and any point $p$, the affine identity \eqref{eq:star} (proved earlier
for $\{0,1,2,3\}$ but valid for any $4$-tuple, by the same barycentric
argument) reads
\[
  [a,b,c,p] \;-\; [a,b,d,p] \;+\; [a,c,d,p] \;-\; [b,c,d,p] \;=\; [a,b,c,d].
\]
Taking signs and noting the four magnitudes on the LHS are positive, the LHS
is forced to be strictly positive (regardless of magnitudes) iff the four
signs $(\sigma(abc), -\sigma(abd), \sigma(acd), -\sigma(bcd))$ are all $+$,
i.e.\ the cell sign pattern at the four ``faces of tetrahedron $abcd$'' is
\[
\bigl(\sigma(abc),\sigma(abd),\sigma(acd),\sigma(bcd)\bigr) \;=\; (+,-,+,-).
\]
Symmetrically the LHS is forced strictly negative iff the pattern is
$(-,+,-,+)$. So:

\noindent\textbf{Forbidden-pattern rule (per $4$-tuple).} A sign vector $\sigma$
violates the chirotope at $\{a,b,c,d\}$ iff
\[
\bigl(\sigma(abc),\sigma(abd),\sigma(acd),\sigma(bcd)\bigr) = \bigl(-,+,-,+\bigr)\cdot \mathrm{sgn}[a,b,c,d],
\]
where the right-hand side means: when $\mathrm{sgn}[a,b,c,d]=+$ the forbidden
pattern is $(-,+,-,+)$, and when $\mathrm{sgn}[a,b,c,d]=-$ the forbidden
pattern is $(+,-,+,-)$. This is exactly the ``missing $16^{\mathrm{th}}$ cell''
pattern from the original section, applied to every tetrahedron formed by four
named points.

\paragraph{Tope characterization.} A sign vector $\sigma$ over the planes
$\binom{[N]}{3}$ corresponds to a non-empty open cell of the realized
configuration iff it does not violate the forbidden-pattern rule at any
$4$-subset $\{a,b,c,d\}$ of named indices. (For a representable rank-$4$
chirotope --- which we have, by construction --- this is necessary and
sufficient: it is the standard Grassmann--Pl\"ucker characterization of topes.)

\paragraph{Slice test.} Given an old cell $C$ with sign vector $\sigma_C$ and a
new plane $\Pi = ij(n{+}1)$, the slice question becomes a small constraint
satisfaction problem. Try to extend $\sigma_C$ to a full sign vector $\sigma$ on
$\binom{\{0,\ldots,n+1\}}{3}$ by assigning a sign to each of the
$\binom{n+1}{2}$ new planes $\bigl\{ab(n{+}1) : 0\le a<b\le n\bigr\}$, subject
to:
\begin{enumerate}
  \item the pinned sign $\sigma(\Pi) = +$;
  \item the forbidden-pattern rule for every $4$-subset
  $\{a,b,c,n{+}1\}$ with $\{a,b,c\}\subset\{0,\ldots,n\}$ (there are
  $\binom{n+1}{3}$ such subsets); each such constraint involves the known
  old-plane sign $\sigma_C(abc)$, the three new-plane signs
  $\sigma(ab(n{+}1)), \sigma(ac(n{+}1)), \sigma(bc(n{+}1))$, and the chirotope
  value $\mathrm{sgn}[a,b,c,n{+}1]$ (also known --- it was pinned when $P_{n+1}$
  was placed in cell $C^\ast$).
\end{enumerate}
$4$-tuples not containing $n{+}1$ are already satisfied by $\sigma_C$ from the
previous step, so they impose nothing new. If at least one extension exists,
the half $C\cap\Pi^+$ is non-empty. Repeat with $\sigma(\Pi) = -$. The plane
slices $C$ iff both halves are non-empty.

\paragraph{Why this works.} Each forbidden-pattern rule is a single
``XOR''-style clause: it forbids exactly one assignment of the three new-plane
signs (the unique completion that, together with $\sigma_C(abc)$ and the
chirotope sign, would form the missing-$16$th-cell pattern of tetrahedron
$\{a,b,c,n{+}1\}$). Out of $2^3 = 8$ possible sign triples on
$\{ab(n{+}1),ac(n{+}1),bc(n{+}1)\}$, exactly $7$ are admissible per clause,
and the slice test is the conjunction over the $\binom{n+1}{3}$ clauses
indexed by the choice of $\{a,b,c\}\subset\{0,\ldots,n\}$.

\paragraph{Why this is also sufficient (and not merely necessary).} Necessity
is immediate from $(\star_{abcd})$: any realized $p$ trivially satisfies every
GP rule. The non-trivial direction is sufficiency --- if $\sigma$ violates no
forbidden pattern, why must some $p$ realize it?

The cleanest justification is via Farkas' lemma. The cell $C\cap\Pi^+$ is
non-empty iff the strict-linear system
\[
\bigl\{\,\sigma(abc)\,[a,b,c,p] > 0 : \text{all triples}\,\bigr\}
\]
is feasible in $p\in\mathbb{R}^3$. By Farkas, infeasibility requires a
non-negative combination $\sum \lambda_{abc}\,\sigma(abc)\,[a,b,c,p]$ that is
identically $\le 0$ as a function of $p$. Since $[a,b,c,p]$ is affine, this
means the linear part vanishes and the constant part is $\le 0$. The
linear-part-vanishing condition is a non-negative-coefficient signed
dependency among the plane-normals $\mathbf{n}_{abc}$.

Now: \emph{every} such dependency comes (up to scaling) from the affine
identity itself. Differentiating $(\star_{abcd})$ in $p$ gives
\[
\mathbf{n}_{abc} \;-\; \mathbf{n}_{abd} \;+\; \mathbf{n}_{acd} \;-\; \mathbf{n}_{bcd} \;=\; 0,
\]
which is the unique alternating $4$-term dependency of normals associated to
the tetrahedron $\{a,b,c,d\}$. The non-negative Farkas combination matching
this dependency forces the sign pattern $\sigma$ to align as $(+,-,+,-)$ or
$(-,+,-,+)$ on those four planes, and then the constant value of the
combination is $\pm[a,b,c,d]$. The ``$\le 0$'' condition on that constant is
\emph{exactly} the GP forbidden pattern.

So: Farkas-infeasibility $\iff$ some GP rule is violated. Contrapositive:
GP-consistent $\iff$ realized. The algorithm is correct in the general case.

(Equivalently, in oriented-matroid language: we are dealing with a
representable rank-$4$ oriented matroid, and Las Vergnas' tope characterization
says topes are exactly the sign vectors satisfying the circuit
sign-conditions, which for rank $4$ are the $(\star_{abcd})$ rules. Mn\"ev's
hardness applies to \emph{abstract} chirotopes; we never face it, because we
extend a realized configuration at every step.)

This is a CSP with $\binom{n+1}{2}$ Boolean variables and $\binom{n+1}{3}$
clauses --- trivially solvable by backtracking for the values of $n$ relevant
here ($n\le 6$ for everything up to the Cs\'asz\'ar polyhedron, giving
$21$ variables and $35$ clauses).

\paragraph{Why this is not NP-complete in our setting.} Realizability of an
arbitrary abstract chirotope is $\exists\mathbb{R}$-complete (Mn\"ev), but that
is not the question we ask: by construction we extend a chirotope that is
already realized. For an already-realized rank-$4$ chirotope, the
Grassmann--Pl\"ucker forbidden-pattern rule \emph{characterizes} topes, so the
slice CSP is a faithful chirotope-only test. The CSP itself has small fixed
arity (each clause touches only three new-plane variables), so it lies far
below the worst-case complexity wall.

\paragraph{Per-step complexity.} For step $n \to n{+}1$: $\binom{n+1}{2}$ new
planes, each tested against every old cell. The number of old cells is
bounded by Zaslavsky's count, $O(n^9)$. Each individual slice test is a CSP
of fixed structure on $\binom{n+1}{2}$ variables and $\binom{n+1}{3}$ clauses,
each of arity $3$; backtracking solves each instance in time polynomial in
$n$. The whole step is therefore polynomial in $n$ (dominated by the cell
count), and entirely chirotope-based.

\end{document}
